% Canonical preamble for this skill. See references/latex-house-style.md.
% The look is references/house-style/ (style-spec.md, housestyle.sty); this
% file adds only the teaching affordances and the math kit on top of it.
% Replace <COURSE>, <DOC>, <NAME> and <YEAR> as needed.
% It ends before \begin{document}, so a smoke document can \input it as-is.
% Build with scripts/build.sh (lualatex, three passes); it puts housestyle.sty
% on the search path and the .sty finds the vendored fonts itself.

\documentclass[11pt]{article}
% A4: uncomment the next line; the margins stay the spec's.
% \PassOptionsToPackage{a4paper}{geometry}
\usepackage{housestyle}  % geometry, fonts, colour tokens, running head, the twelve blocks, hyperref (hidelinks)
\usepackage{amssymb,amsthm}
\usepackage{mathtools}
\usepackage{tabularx,multirow}

% --- Running head and foot -------------------------------------------------
% Slug at the left of the head; the section name at the right follows each
% \section (override with \hssection{...}). Foot: the copyright line at the
% left, the page number at the right. The first page has no head, per the spec.
\hsslug{<COURSE> <DOC>}
\fancyfoot[L]{\hslabel{\textcopyright{} <YEAR> <NAME>. All rights reserved.}}
\makeatletter
\renewcommand{\sectionmark}[1]{\gdef\@hssection{#1}}
\makeatother

% --- Headings --------------------------------------------------------------
% Headings are unnumbered on the page (the .sty's \section). Cross-reference
% a section by name with \nameref{sec:...}; a Part by its number, Part~\ref{part:...}.
% A Part opens on a new page: its eyebrow, then the page heading at 28/34, no bold.
\makeatletter
\renewcommand{\thepart}{\arabic{part}}
\newcommand{\hs@parteyebrow}{}
\titleclass{\part}{top}
\assignpagestyle{\part}{fancy}  % keep the head and the copyright foot on a Part page
\titleformat{\part}[display]{\fontsize{28pt}{34pt}\selectfont\color{ink}}%
  {\hsaccentlabel{Part \thepart\hs@parteyebrow}}{10pt}{}
\titlespacing*{\part}{0pt}{0pt}{24pt}
\titleformat{\subsection}{\fontsize{14pt}{19pt}\selectfont\color{ink}}{}{0pt}{}
\titlespacing*{\subsection}{0pt}{18pt}{6pt}
\titleformat{\subsubsection}{\itshape\color{ink}}{}{0pt}{}
\titlespacing*{\subsubsection}{0pt}{12pt}{2pt}

% --- Taxonomy: category as structure, never as colour ------------------------
% \catbanner{A}{Title} opens a Part: eyebrow "Part N . Category A", then the
% heading, and the running head's section becomes the title.
\newcommand{\catbanner}[2]{%
  \gdef\hs@parteyebrow{\ \textperiodcentered\ Category #1}%
  \part{#2}\hssection{#2}\gdef\hs@parteyebrow{}}
\makeatother
\newcommand{\cat}[1]{\hslabel{[#1]}}      % [A]
\newcommand{\tool}[1]{\hslabel{[T#1]}}    % [T5]
\newcommand{\tools}[1]{\hslabel{[#1]}}    % [T4,T5]

% --- Teaching affordances, inside the twelve blocks -------------------------
% insight: the section's lead line, right under the heading.
\newcommand{\insight}[1]{\hslead{#1}}
% derivation: a display-equation run (block 12) the reader can skip on
% re-read. A label opens it, a grey rule shows where it ends.
\newenvironment{derivation}%
  {\par\vspace{4pt}\hslabel{Derivation \textperiodcentered\ skip on re-read}\par\nopagebreak}%
  {\par\nopagebreak\vspace{2pt}\noindent\hsthinrule\par}
% \onsheet: this equation lives on the formula sheet. Grey, not accent.
\newcommand{\onsheet}{\par\nopagebreak\noindent\hslabel{On the formula sheet}\nopagebreak}
% \opt: out of scope. Inline after a mention, or first thing in an optional paragraph.
\newcommand{\opt}{\leavevmode\unskip\hspace{4pt}\hslabel{[optional]}\hspace{4pt}\ignorespaces}
% \probhead{source}{method}{category}: one mono label line over a worked problem.
\newcommand{\probhead}[3]{%
  \par\vspace{8pt}\noindent\hslabel{#1 \textperiodcentered\ #2 \textperiodcentered\ #3}\par\nopagebreak}

% --- Tables ----------------------------------------------------------------
% L{w} and R{w} come from the .sty; Y is a ragged-right tabularx column.
\newcolumntype{Y}{>{\raggedright\arraybackslash}X}

% --- Math kit ----------------------------------------------------------------
\newcommand{\R}{\mathbb{R}}
\newcommand{\N}{\mathbb{N}}
\newcommand{\Z}{\mathbb{Z}}
\newcommand{\RR}[1]{\mathbb{R}^{#1}}
\newcommand{\grad}{\nabla}
\newcommand{\Hess}{H_f}

% Bold vectors -- define all lowercase so agents can grab any letter
\newcommand{\ba}{\mathbf{a}}  \newcommand{\bb}{\mathbf{b}}  \newcommand{\bc}{\mathbf{c}}
\newcommand{\bd}{\mathbf{d}}  \newcommand{\be}{\mathbf{e}}  \newcommand{\bff}{\mathbf{f}}
\newcommand{\bg}{\mathbf{g}}  \newcommand{\bh}{\mathbf{h}}  \newcommand{\bi}{\mathbf{i}}
\newcommand{\bj}{\mathbf{j}}  \newcommand{\bk}{\mathbf{k}}  \newcommand{\bl}{\mathbf{l}}
\newcommand{\bm}{\mathbf{m}}  \newcommand{\bn}{\mathbf{n}}  \newcommand{\bp}{\mathbf{p}}
\newcommand{\bq}{\mathbf{q}}  \newcommand{\br}{\mathbf{r}}  \newcommand{\bs}{\mathbf{s}}
\newcommand{\bt}{\mathbf{t}}  \newcommand{\bu}{\mathbf{u}}  \newcommand{\bv}{\mathbf{v}}
\newcommand{\bw}{\mathbf{w}}  \newcommand{\bx}{\mathbf{x}}  \newcommand{\by}{\mathbf{y}}
\newcommand{\bz}{\mathbf{z}}  \newcommand{\bzero}{\mathbf{0}}

% Transpose: braces (not ^) so X^\trans gives X^{\!\top}
\newcommand{\trans}{{\!\top}}

% Partial derivatives
\newcommand{\pderiv}[2]{\dfrac{\partial #1}{\partial #2}}
\newcommand{\ppderiv}[2]{\dfrac{\partial^2 #1}{\partial #2^2}}
\newcommand{\pmderiv}[3]{\dfrac{\partial^2 #1}{\partial #2 \, \partial #3}}

% Big-O -- use \Oh{n} or \Ohof{n \log n} everywhere; never bare $O()$
\newcommand{\Oh}{\mathcal{O}}
\newcommand{\Ohof}[1]{\mathcal{O}\!\left(#1\right)}

% Differential d
\newcommand{\diff}[1]{\mathop{}\!\mathrm{d}#1}

% --- Code --------------------------------------------------------------------
% Block 11: \hslisting{Listing 1 / what it shows} then
% \begin{Verbatim}[bgcolor=codefill] ... \end{Verbatim}. No syntax colour in print.

% \hsnumbersections  % uncomment for numbered sections (1, 1.1, ...)
\hsslug{<Course> <Doc>}
\hsdate{<Month Year>}
\begin{document}
\begin{titlepage}
\thispagestyle{hstitlepage}
\vspace*{\fill}
\hsaccentlabel{<Course> \textperiodcentered\ Document type}\par\vspace{14pt}
\hstitle[44pt]{Document title goes here}\vspace{20pt}
\hslead{A one-sentence lead that says what this document is for and who should read it.}
\vspace*{\fill}\vspace*{\fill}
\end{titlepage}
\setcounter{page}{2}
\tableofcontents
\clearpage

\catbanner{A}{First part title}
\section{First section}
\hslead{A lead paragraph in the larger grey face, opening the section.}

Body text sits here. Lorem ipsum dolor sit amet, consectetur adipiscing elit,
sed do eiusmod tempor incididunt ut labore et dolore magna aliqua. Ut enim ad
minim veniam, quis nostrud exercitation ullamco laboris nisi ut aliquip ex ea
commodo consequat.

\subsection{A subsection}
Placeholder paragraph with an inline equation $f(\bx) = \bx^\trans A \bx$ and a
displayed one:
\begin{equation}
\grad f(\bx) = (A + A^\trans)\,\bx.
\end{equation}
Here $A$ is the matrix of the quadratic form and $\bx$ the point at which the
gradient is taken.

\begin{itemize}
\item First bullet point.
\item Second bullet point.
\end{itemize}

\begin{hsstatrow}[3]
\hsstat{42}{Caption for the first figure}
\hsstat{7.5}{Caption for the second figure}
\hsstat{100\%}{Caption for the third figure}
\end{hsstatrow}

\begin{hscallout}{Callout label}
A callout holds one point the reader must not miss. Placeholder text.
\end{hscallout}

\hsclaim{A single pull-quote claim, set large between two rules.}{Attribution label}

\clearpage
\section{How the pieces connect}
\begin{hsfigure}{Figure 1 / node-and-arrow flow}%
{Slate boxes are inputs and outputs. Sage boxes are work an agent does.
Rust boxes are deterministic checks that pass or fail.}
\begin{tikzpicture}[node distance=14pt]
  \node[hsnode] (a) {\hsnodetext{Input}{A brief}{goal, boundaries, done-criteria}};
  \node[hswork,right=of a] (b) {\hsnodetext{Agent}{Draft}{structure first, then prose}};
  \node[hsgate,right=of b] (c) {\hsnodetext{Gate}{Style check}{pass, or back with reasons}};
  \node[hsnode,right=of c] (d) {\hsnodetext{Output}{A PDF}{with its sources attached}};
  \draw[hsarrow] (a) -- (b); \draw[hsarrow] (b) -- (c); \draw[hsarrow] (c) -- (d);
  \draw[hsfail] (c.south) -- ++(0,-12pt) -| (b.south);
\end{tikzpicture}
\hsfailnote{A failed gate returns to the draft step, twice, then to a person.}
\end{hsfigure}

\hsprovenance{Figure 1 from core/docs/DESIGN.md at commit 4f2a9c1, 21 September 2026.}
% Two diagrams do not fit one page with their legends; a figure that breaks
% leaves its legend alone on the next page.
\clearpage
\begin{hsfigure}{Figure 2 / what runs with nobody there}%
{Slate boxes are doors into the box, the rust one is what every event goes
through, the sage one is where a model runs and the grey one is a record. An
italic line on an arrow says what passes between two parts.}
\begin{tikzpicture}[x=1pt,y=1pt]
  \node[hsrole,hsslate]      (d) at (0,72)    {\hsrolename{Slack daemon}{the box's one connection}};
  \node[hsrole,hsslate]      (m) at (0,0)     {\hsrolename{Mail receiver}{loopback, behind a tunnel}};
  \node[hsrole,hsaccentrole] (b) at (170,36)  {\hsrolename{Broker}{every 60 s: the event door}};
  \node[hsrole,hssage]       (s) at (340,72)  {\hsrolename{Sessions}{the tmux windows}};
  \node[hsrole,hsquiet]      (r) at (170,-52) {\hsrolename{The board}{one row per task}};
  \draw[hsarrow] (d) -- node[hsann] {stored,\\then handed on} (b);
  \draw[hsarrow] (m) -- (b);
  \draw[hsarrow] (b) -- node[hsann] {one message,\\not six} (s);
  \draw[hsarrow] (b) -- (r);
\end{tikzpicture}
\end{hsfigure}

\hsprovenance{Figure 2 from core/docs/DESIGN.md at commit 4f2a9c1, 21 September 2026.}
\clearpage
\section{Second section}
\begin{hsblock}
\hslabel{Table 1 / Placeholder table}\par\vspace{8pt}
\begin{tabular}{@{}lll@{}}
\hstoprule
\hshead{Column A} & \hshead{Column B} & \hshead{Column C} \\
\hline
Row one & value & value \\
Row two & value & value \\
\hline
\end{tabular}
\end{hsblock}

\hslisting{Listing 1 / Placeholder code}
\begin{Verbatim}[bgcolor=codefill]
def placeholder(x):
    return x * 2
\end{Verbatim}

% Sources open their own page: the running head takes the last heading on a
% page, so a Sources heading lower down would retitle the page above it.
\clearpage
\section*{Sources}
\begin{hssources}
\item Author, \emph{Title of the first source}, Publisher, Year.
\item Author, \emph{Title of the second source}. \hsurl{https://example.com/source}
\end{hssources}
\end{document}
